\documentclass[11pt,letterpaper]{article}
\usepackage[T1]{fontenc}
\usepackage[utf8]{inputenc}
\usepackage{lmodern,microtype}
\usepackage[margin=.92in,headheight=15pt]{geometry}
\usepackage{amsmath,amssymb,amsthm,mathtools}
\usepackage{booktabs,tabularx,array,longtable,enumitem}
\usepackage{xcolor,fancyhdr,listings,xurl,needspace}
\usepackage[colorlinks=true,linkcolor=blue!45!black,urlcolor=blue!45!black,citecolor=blue!45!black]{hyperref}
\definecolor{ink}{RGB}{24,47,65}
\definecolor{pale}{RGB}{245,248,250}
\setlength{\emergencystretch}{3em}
\setlength{\parskip}{.3em}
\setlist{itemsep=3pt,topsep=5pt,leftmargin=1.7em}
\newtheorem{theorem}{Theorem}[section]
\newtheorem{lemma}[theorem]{Lemma}
\newtheorem{proposition}[theorem]{Proposition}
\newtheorem{corollary}[theorem]{Corollary}
\theoremstyle{definition}\newtheorem{definition}[theorem]{Definition}
\theoremstyle{remark}\newtheorem{remark}[theorem]{Remark}
\newcommand{\ipra}{\mathsf{iPRA}}
\newcommand{\TJ}{\mathsf T_J}
\newcommand{\TI}{\mathsf T_{\mathrm{eq}}}
\newcommand{\FV}{\operatorname{FV}}
\newcommand{\J}[2]{J(#1;#2)}
\newcommand{\biff}{\mathrel{\leftrightarrow}}
\newcommand{\file}[1]{{\small\nolinkurl{#1}}}
\lstset{basicstyle=\ttfamily\footnotesize,columns=fullflexible,keepspaces=true,
 breaklines=true,showstringspaces=false,frame=single,rulecolor=\color{black!20},
 backgroundcolor=\color{pale},aboveskip=8pt,belowskip=8pt}
\pagestyle{fancy}\fancyhf{}
\fancyhead[L]{\small A register-graph presentation of iPRA}
\fancyhead[R]{\small November 10, 2026}
\fancyfoot[C]{\thepage}
\hypersetup{pdftitle={A Register-Graph Presentation of iPRA},
 pdfauthor={Research report on the ConciseIPRA proposal}}
\title{\color{ink}\bfseries A Register-Graph Presentation of iPRA\\[8pt]
\large Removing primitive successor and initialization\\with an explicit whole-domain bi-interpretation}
\author{Research report on the ConciseIPRA proposal}
\date{November 10, 2026}
\begin{document}
\maketitle
\begin{abstract}
We replace the numerical syntax $0,S,I(n,a;u.h)$ of the current
\emph{ConciseIPRA} proposal by $0,J(n;i,r.h)$. The register of $J$ always
starts at one. A zero instruction increments it; an instruction $S(t)$
loads $t$. Successor is the abbreviation $S(a)=J(a;i,r.0)$, and no
independent initializer is primitive. A transition-graph biconditional
also supplies numerical coverage. Predecessor and an unrestricted
selector have short derivations without induction. Every primitive
recursion then compiles by one extra iteration, again without induction
or pairing. Literal-equation induction recovers quantifier-free
induction, and both interpretation composites are internally equal to
the identity on the entire numerical domain.

This is a mathematical simplification of the term grammar and arithmetic
bootstrap. It is \emph{not} an established improvement on the smallest
complete proof-checker size: the transition axiom is richer, and the
counter-aware body has two bound variables rather than one. We supply a
readable complete checker for the new theory and reproducible measurements.
A separate 288-AST-node implementation of the inherited scalar calculus
is included as an ancillary implementation result; it does not implement
the new arithmetic. No absolute simplest axiomatization or
proof-assistant verification is claimed.
\end{abstract}
\begin{center}
\fcolorbox{ink}{pale}{\begin{minipage}{.91\linewidth}
\textbf{Main construction.}
\[
 \boxed{\quad t::=x\mid0\mid J(n;i,r.h),\qquad S(a):=J(a;i,r.0).\quad}
\]
The initializer is recovered by the explicit compiler
\[
 R(n,a;u,v.h)=J\bigl(Sn;i,r.S(C(i,a,h(Pi,r)))\bigr).
\]
Its equations are $R(0,a;h)=a$ and
$R(Sn,a;h)=h(n,R(n,a;h))$.
\end{minipage}}
\end{center}
\medskip
\noindent\textbf{Reading guide.} The complete theory is in
Section~\ref{sec:theory}; the direct recursion compiler is in
Section~\ref{sec:compiler}; equality reflection and the two internal
round trips are in Sections~\ref{sec:reflection}--\ref{sec:interpretations}.
Section~\ref{sec:artifacts} separates the mathematical result from
the executable-checker measurements.
\newpage

\section{Result and comparison boundary}

The comparison is pinned to ProveIt commit
\texttt{8662bc0a00f4e77e482f7d141d7a4f919f70f9aa}, read on
November 10, 2026 in the user's time zone. Its
\file{Logic/PrimitiveRecursiveArithmetic/ConciseIPRA/README.md}
describes a native scalar calculus with terms $0,S,I$, formulas built
from equality and two binders, and an ordinary realization $\TI$.
The ordinary realization has intuitionistic equality logic, scalar
iterator equations, numerical coverage, $0\ne S0$, and contextual
literal-equation induction \cite{repo,basis}.

The new arithmetic uses ordinary first-order intuitionistic logic.
The logical rules are specified in Section~\ref{sec:logic}; they are
part of the presentation and are included in the executable checker.
We do not compare an arithmetic-only program against a complete old
checker.

\begin{theorem}[Main result]\label{thm:main}
The theory $\TJ$ defined in Section~\ref{sec:theory} is
bi-interpretable with full first-order intuitionistic primitive
recursive arithmetic with quantifier-free induction. Both domains are
the whole numerical domain, interpreted equality is actual equality,
and the graph of each composite isomorphism is $x=y$.
Consequently the two theories are equiconsistent.
\end{theorem}

Here ``full first-order'' concerns the available formulas and logical
rules. Induction is restricted to quantifier-free formulas in $\ipra$,
and initially to literal equations in $\TJ$. No induction on arbitrary
quantified formulas is smuggled into the argument.

\begin{center}\small
\begin{tabularx}{\linewidth}{@{}lXX@{}}
\toprule
Feature & Current scalar arithmetic & Proposed register arithmetic\\
\midrule
Primitive numerical constructors & $0,S,I$ & $0,J$\\
Independent initializer & Argument of $I$ & Derived by one extra step\\
Local variables in a body & One: previous state & Two: old counter and state\\
Predecessor bootstrap & A scalar iteration construction & Four short abbreviations\\
General primitive recursion & Derived using pairing & Direct one-line compiler\\
Computation specification & Two equations & One base equation and a graph biconditional\\
Numerical coverage & Separate in $\TI$ & Derived from the graph\\
Complete-checker byte record & 529 normalized bytes & Not improved here\\
\bottomrule
\end{tabularx}
\end{center}

The counter variable is a real cost, as is the richer graph axiom.
A one-body-variable version is given in Appendix~\ref{sec:scalar}.
That version retains the old binding arity while still eliminating
primitive successor and the initializer, but its simulation of general
primitive recursion passes through the scalar bridge.

\section{Complete theory}\label{sec:theory}

\subsection{Numerical syntax and finite descriptions}

Terms are finite expressions
\[
 t::=x\mid0\mid J(n;i,r.h).
\]
The distinct variables $i,r$ are bound \emph{only in $h$}; they are not
bound in $n$. Other variables in $h$ are passive numerical parameters.
Substitution is capture-free. We use the abbreviations
\[
 S(a):=J(a;i,r.0),\qquad 1:=S0,\qquad
 f(n,a):=J(n;i,r.a),
\]
where the bound variables are chosen fresh for $a$. Fixed superscripts
$S^k,P^k$ mean finite composition, never exponentiation of an unspecified
function.

The binding notation is a finite-description notation for a first-order
language, not an additional function sort. More formally, start with
the fixed zero-body description $J_0$. For every fixed finite body
$h(i,r,\vec p)$ built from earlier descriptions, introduce the
first-order function symbol $J_h(n,\vec p)$ and its displayed schemas.
The symbol $S$ is always the abbreviation $J_0$.
Nested bodies are handled in the same finite dependency order.
The zero-body schemas themselves use $S=J_0$; their standard
interpretation is established in Section~\ref{sec:interpretations}.
No variable ranges over descriptions and no uniform evaluator is
postulated inside the theory.

\subsection{Logical language and rules}\label{sec:logic}

Formulas are generated by
\[
 A::=(s=t)\mid\bot\mid(A\land A)\mid(A\lor A)\mid(A\to A)
       \mid\forall x A\mid\exists x A.
\]
Negation abbreviates $A\to\bot$, and $A\leftrightarrow B$
abbreviates $(A\to B)\land(B\to A)$.
Contexts $\Gamma$ are finite reusable collections of assumptions.
The logical inference rules are the following ordinary natural-deduction
rules. Displayed auxiliary assumptions are discharged.

\begin{center}\small
\begin{tabularx}{\linewidth}{@{}lXX@{}}
\toprule
Rule & Premises & Conclusion\\
\midrule
Hypothesis & $A\in\Gamma$ & $\Gamma\vdash A$\\
Equality reflexivity & none & $\Gamma\vdash t=t$\\
Equality transport & $\Gamma\vdash s=t$, $\Gamma\vdash F(s)$ & $\Gamma\vdash F(t)$\\
Falsity elimination & $\Gamma\vdash\bot$ & $\Gamma\vdash A$\\
$\land$ introduction & $\Gamma\vdash A$, $\Gamma\vdash B$ & $\Gamma\vdash A\land B$\\
$\land$ elimination & $\Gamma\vdash A\land B$ & Either conjunct\\
$\lor$ introduction & $\Gamma\vdash A$ or $\Gamma\vdash B$ & $\Gamma\vdash A\lor B$\\
$\lor$ elimination & $\Gamma\vdash A\lor B$, $\Gamma,A\vdash F$, $\Gamma,B\vdash F$ & $\Gamma\vdash F$\\
$\to$ introduction & $\Gamma,A\vdash B$ & $\Gamma\vdash A\to B$\\
$\to$ elimination & $\Gamma\vdash A\to B$, $\Gamma\vdash A$ & $\Gamma\vdash B$\\
$\forall$ introduction & $\Gamma\vdash A(z)$, $z$ fresh & $\Gamma\vdash\forall x A(x)$\\
$\forall$ elimination & $\Gamma\vdash\forall x A(x)$ & $\Gamma\vdash A(t)$\\
$\exists$ introduction & $\Gamma\vdash A(t)$ & $\Gamma\vdash\exists x A(x)$\\
$\exists$ elimination & $\Gamma\vdash\exists x A(x)$, $\Gamma,A(z)\vdash F$, $z$ fresh & $\Gamma\vdash F$\\
\bottomrule
\end{tabularx}
\end{center}

For universal introduction, $z$ is absent from $\Gamma$ and the fixed
parameters of $A$. For existential elimination, it is absent from
$\Gamma,F$ and the fixed parameters of the matrix. The template $F(v)$
in equality transport is any formula. Weakening, exchange, contraction,
renaming, and assumption cut are admissible. There is no classical rule
or excluded-middle axiom. These rules, together with the arithmetic
schemas below, are the complete presentation.

\subsection{Arithmetic schemas}

For each fixed body $h(i,r,\vec p)$, write $j_n=J(n;i,r.h)$ and
$e_n=h(n,j_n,\vec p)$. Use universal closure over the free numerical
variables in the following schemas:
\begin{align}
 J(0;i,r.h)&=1, &&\text{(J0)}\label{ax:J0}\\
 J(Sn;i,r.h)=b&\ \leftrightarrow\
       \bigl((e_n=0\land b=Sj_n)\lor e_n=Sb\bigr),
       &&\text{(STEP)}\label{ax:step}\\
 0=1&\ \to\bot. &&\text{(Sep)}\label{ax:sep}
\end{align}
In STEP, $b$ is an output variable distinct from the body's passive
parameters before universal instantiation. All instances obtained by
capture-free substitution are allowed.

The one arithmetic inference rule is contextual equation induction:
\[
 \frac{\Gamma\vdash s(0,\vec p)=t(0,\vec p)\qquad
       \Gamma,s(z,\vec p)=t(z,\vec p)
                \vdash s(Sz,\vec p)=t(Sz,\vec p)}
      {\Gamma\vdash s(a,\vec p)=t(a,\vec p)}.
 \tag{EI}\label{rule:EI}
\]
The variable $z$ is fresh for $\Gamma$, the target $a$, and the fixed
parameters. The invariant must be a \emph{literal equation}; an equation
under a quantifier is not eligible. Contexts may contain arbitrary
first-order formulas. No further arithmetic axiom is primitive.

\subsection{Operational meaning and intrinsic coverage}

In the standard interpretation, $J(n;h)$ executes $n$ instructions with
an old-counter variable $i$ and a register $r$, initially $r=1$.
The fixed finite body returns an instruction value. Zero increments
the register; a value $S(t)$ loads $t$. Thus, for each fixed body,
this is ordinary primitive recursion.

The right-to-left direction of STEP immediately gives the two useful
conditional computation rules
\begin{align}
 e_n=0&\ \Longrightarrow\ J(Sn;h)=Sj_n,\label{rule:inc}\\
 e_n=St&\ \Longrightarrow\ J(Sn;h)=t.\label{rule:load}
\end{align}
Their premises are proved equations, not tests performed by an
untrusted evaluator.

\begin{lemma}[Functional coverage]\label{lem:coverage}
Define $W(a):=f(1,a)$. Before any induction,
\[
 \TJ\vdash a=0\lor a=S(W(a)).
\]
In particular, $\forall a(a=0\lor\exists x\,a=Sx)$ is derivable.
\end{lemma}
\begin{proof}
Take the constant body $h(i,r)=a$, set $n=0$ and $b=f(1,a)$ in
STEP, and use reflexivity on its left side. Its right side is
$(a=0\land b=S1)\lor a=Sb$. Forget the unused conjunct.
\end{proof}

Conversely, J0, the two conditional rules, and numerical coverage
prove STEP. For its left-to-right direction, split coverage for $e_n$.
At zero use the increment equation; at $St$ use the load equation
and the given equality of outputs. The right-to-left direction is
computation. This equivalence does not require induction, separation,
or successor injectivity. The biconditional is essential: its
left-to-right direction alone would not assert deterministic computation.

One can package J0 and STEP into a single graph schema by explicitly
distinguishing a zero counter from a successor counter. Such packaging
also works for the old iterator, so a smaller number of named schemas
alone is not the claimed improvement.

\section{A short arithmetic bootstrap}

All endpoint equations in this section use only equality and the
right-to-left direction of STEP together with J0. They need neither
induction nor coverage. Consequently none of the constructions relies
on already having arithmetic comparison or an equality characteristic.

\subsection{Constant instructions and predecessor}

The abbreviation $f(n,a)=J(n;i,r.a)$ satisfies
\begin{equation}
 f(0,a)=1,\qquad f(n,0)=Sn,\qquad f(Sn,Sa)=a.
 \label{eq:f}
\end{equation}
The middle equation is the definition of $S$. The last is a load
step whose premise is reflexivity $Sa=Sa$.

Define
\[
 Z(x)=f(x,1),\quad W(x)=f(1,x),\quad
 D(x)=f(x,Sx),\quad P(x)=W(D(x)).
\]
Then direct instances of \eqref{eq:f} give
\begin{alignat}{3}
 Z(0)&=1,&\qquad Z(Sx)&=0,\label{eq:Z}\\
 W(0)&=S1,& W(Sx)&=x,\label{eq:W}\\
 D(0)&=1,& D(Sx)&=Sx,\label{eq:D}\\
 P(0)&=0,& P(Sx)&=x.\label{eq:P}
\end{alignat}
For example, $P(0)=W(1)=0$ and $P(Sx)=W(Sx)=x$.
Thus $P$ is a predecessor throughout the internal domain, and successor
injectivity follows by applying $P$ to $Sx=Sy$.

Using Sep, successor is never zero: an equality $Sx=0$ implies
$Z(Sx)=Z0$, hence $0=1$, and therefore $\bot$.
These facts are consequences, not successor axioms silently appended
to the presentation.

\subsection{Unrestricted selection}

Put
\[
 G(c,a)=P^2 f(Sa,Zc).
\]
Because $Z(0)=1$ and $Z(Sn)=0$, equations \eqref{eq:f}
and \eqref{eq:P} give
\begin{equation}
 G(0,a)=0,\qquad G(Sn,a)=a.\label{eq:G}
\end{equation}
Indeed, the inner $f$ is zero in the first case and $S^2a$ in the
second case. Now define
\begin{equation}
 C(c,a,b):=P^2 f\bigl(Sa,\,G(c,S^3b)\bigr).
 \label{def:C}
\end{equation}
At $c=0$, the instruction $G(c,S^3b)$ is zero, so the inner $f$
is $S^2a$. At $c=Sn$, the instruction is $S^3b$, and the
positive counter $Sa$ loads $S^2b$. Hence
\begin{equation}
 \boxed{\ C(0,a,b)=a,\qquad C(Sn,a,b)=b.\ }\label{eq:C}
\end{equation}
This is a selector of arbitrary numerical outputs. It is not limited
to Boolean results or to values bounded by a preselected initializer.
The direct use of the old register in the increment branch is what
escapes the single-body information restriction of tied-seed
iterations discussed in the current research packages \cite{barrier}.

\section{Every primitive recursion in one extra step}\label{sec:compiler}

Consider any fixed intended recursion with passive parameters:
\[
 F(0,\vec p)=a(\vec p),\qquad
 F(Sn,\vec p)=h(n,F(n,\vec p),\vec p).
\]
Assume $a,h$ have already been translated to $J$ terms. Choose fresh
local variables $i,r$ and define
\begin{align}
 H(i,r,\vec p)&:=S\bigl(C(i,a(\vec p),h(Pi,r,\vec p))\bigr),\label{def:H}\\
 R(n,a;u,v.h)&:=J(Sn;i,r.H(i,r,\vec p)).\label{def:R}
\end{align}
The old counter distinguishes the first update from the subsequent ones.
The instruction is always a successor, so every update is a load.

\begin{theorem}[Direct recursion compiler]\label{thm:compiler}
Without induction, coverage, pairing, or equality reflection,
\[
 R(0,a;h)=a,\qquad R(Sn,a;h)=h(n,R(n,a;h)).
\]
\end{theorem}
\begin{proof}
The first update has counter $i=0$ and old register $1$. By
\eqref{eq:C}, its instruction is $Sa$, and one load gives
$J(1;H)=a$.

For the next equation, put $r=J(Sn;H)=R(n,a;h)$. Its old counter
is $Sn$, so \eqref{eq:C} and $P(Sn)=n$ give
\[
 H(Sn,r)=S\bigl(h(n,r)\bigr).
\]
One load therefore gives
$J(S(Sn);H)=h(n,J(Sn;H))$, exactly the required equation.
\end{proof}

Initial functions and composition translate by literal term substitution;
Theorem~\ref{thm:compiler} handles each primitive-recursion node.
External induction on the finite description constructs a $J$ term
for every primitive-recursive description and proves both of its
defining equations internally. Nothing in this compiler quantifies
over a class of functions inside the theory.

For example, define
\[
 x+y=R(y,x;u,v.Sv),\qquad
 d(x,y)=R(y,x;u,v.Pv),\qquad
 x\cdot y=R(y,0;u,v.v+x).
\]
Their ordinary recursion equations follow immediately. These are
finite $J$ terms after expansion. A compact source grammar does not
guarantee small expanded terms; the archive's actual syntax compiler
makes that separate cost inspectable.

\section{From equation induction to quantifier-free induction}\label{sec:reflection}

This section adapts the repository's residual-reflection proof
\cite{reflection}, using the already available $Z,W,P,C$. The proof
is included to keep the main equivalence independent of a hidden
equality-test assumption. All internal inductions below have a single
literal equation as invariant.

\subsection{Elementary identities}

The recursion equations for addition and $d$ are
\begin{alignat}{3}
 x+0&=x,&\qquad x+Sy&=S(x+y),\label{eq:add}\\
 d(x,0)&=x,& d(x,Sy)&=P(d(x,y)).\label{eq:monus}
\end{alignat}
Equation induction yields
\begin{align}
 0+y&=y,& (Sx)+y&=S(x+y),\label{eq:addderived}\\
 d(Sx,Sy)&=d(x,y),&d(Sa,a)&=1.\label{eq:dderived}
\end{align}
For the first two identities, induct on $y$ using \eqref{eq:add}.
For cancellation, the zero case is $P(Sx)=x$; the successor step
applies $P$ to the induction equation. For the diagonal, the zero
case is $d(1,0)=1$ and the successor case is cancellation.
In particular, $(Sy)+r=y+Sr$ follows from \eqref{eq:addderived}.

Coverage, the constructor laws, and \eqref{eq:Z} prove
\[
 Z(v)=0\lor Z(v)=1,\qquad Z(v)=1\leftrightarrow v=0.
\]
Likewise \eqref{eq:W} proves
\begin{equation}
 W(v)=0\leftrightarrow v=1.\label{eq:Wtest}
\end{equation}
For the forward implication in \eqref{eq:Wtest}, split coverage of
$v$. Its zero branch gives $W0=2\ne0$, a contradiction; its
successor branch $v=Sx$ gives $x=0$. No induction is involved.

\subsection{One residual equation reflects actual equality}

\begin{lemma}\label{lem:residual}
For every $x,y$,
\begin{equation}
 C\bigl(d(x,y),x,y+d(x,y)\bigr)=x.\label{eq:residual}
\end{equation}
\end{lemma}
\begin{proof}
Hold $x$ fixed and induct on $y$. At zero the equation is
$C(x,x,0+x)=x$; both branches equal $x$, by \eqref{eq:addderived}
and coverage.

For the step write $r=d(x,y)$ and $r'=Pr=d(x,Sy)$.
If $r'=0$, the new selector directly returns $x$.
Otherwise coverage writes $r'$ as a successor. Coverage of $r$
cannot choose zero, since that would make $r'=P0=0$; in its
successor branch $r=S(Pr)=Sr'$. The induction equation now
reduces to $y+r=x$. Meanwhile
\[
 (Sy)+r'=y+Sr'=y+r=x.
\]
This is the branch selected by the successor control $r'$ in
the new equation. Every split is supplied by coverage; the
induction invariant itself remains \eqref{eq:residual}.
\end{proof}

Applying the lemma at $x=Sa,y=b$ shows
\begin{equation}
 d(Sa,b)=1\leftrightarrow a=b.\label{eq:resreflect}
\end{equation}
In the forward direction it gives $b+1=Sa$, hence $Sb=Sa$
and then $b=a$ by $P$. The converse is the diagonal identity.
Thus the finite term
\begin{equation}
 e(a,b):=Z\bigl(W(d(Sa,b))\bigr)\label{def:equality}
\end{equation}
has the internally proved properties
\begin{equation}
 e(a,b)=0\lor e(a,b)=1,\qquad e(a,b)=1\leftrightarrow a=b.
 \label{eq:equality}
\end{equation}
It tests actual equality on the whole domain, not equality only
between standard numerals or a separately postulated congruence.

\subsection{Boolean compilation of formulas}

For each fixed quantifier-free formula $A$, define its characteristic
term by recursion on its finite syntax:
\begin{align*}
 \chi_{s=t}&=e(s,t),&\chi_{\bot}&=0,\\
 \chi_{A\land B}&=C(\chi_A,0,\chi_B),&
 \chi_{A\lor B}&=C(\chi_A,\chi_B,1),\\
 \chi_{A\to B}&=C(\chi_A,1,\chi_B).
\end{align*}
An external structural induction proves inside $\TJ$
\begin{equation}
 (\chi_A=0\lor\chi_A=1)\quad\land\quad
 (\chi_A=1\leftrightarrow A).\label{eq:bool}
\end{equation}
For atoms use \eqref{eq:equality}; for $\bot$ use Sep and explosion.
For conjunction, the zero antecedent branch returns zero and excludes
$A$, while the one branch returns $\chi_B$. For disjunction, those
branches return $\chi_B$ and one, respectively. For implication, the
zero branch proves $\neg A$ and returns one, while the one branch
already proves $A$ and returns $\chi_B$. These are constructive
case arguments using the proved two-value range, not excluded middle
for arbitrary formulas.

Given a quantifier-free base $A(0)$ and a contextual induction step
$A(z)\vdash A(Sz)$, use \eqref{eq:bool} to obtain the corresponding
base and step for the literal equation $\chi_A=1$. EI gives that
equation at any target, and \eqref{eq:bool} recovers $A$ there.
The retained context is unchanged, including any quantified assumptions.
This proves full contextual quantifier-free induction.

\section{The interpretations and both internal round trips}\label{sec:interpretations}

\subsection{Interpreting \texorpdfstring{$\TJ$}{T-J} in \texorpdfstring{$\ipra$}{iPRA}}

Translate each fixed body recursively to primitive-recursive syntax.
Using the standard predecessor and selector of $\ipra$, define
\[
 J_h(0,\vec p)=1,
\]
\[
 J_h(Sn,\vec p)=
 C\bigl(h(n,J_h(n,\vec p),\vec p),\,
       S(J_h(n,\vec p)),\,
       P(h(n,J_h(n,\vec p),\vec p))\bigr).
 \tag{J in iPRA}
\]
This is a valid fixed primitive-recursive description. For the zero
body, equation induction proves $J_0(n)=S_{\ipra}n$: the base
is one and the step increments. Therefore the new abbreviation $S$
is interpreted as the old successor, provably at every internal
number, not just in the standard model.

J0 and Sep now translate to theorems. To prove STEP, use numerical
coverage for the instruction. At zero the update increments, and at
a successor it loads its predecessor. Both directions of the
biconditional follow by equality reasoning. EI translates to a
quantifier-free induction instance. Ordinary logic and equality
are preserved literally.

\subsection{Interpreting \texorpdfstring{$\ipra$}{iPRA} in \texorpdfstring{$\TJ$}{T-J}}

Translate zero and variables identically and translate successor by
$J_0$. Translate initial functions and composition directly, and
each primitive-recursion node using Theorem~\ref{thm:compiler}.
Its defining equations hold internally. Successor injectivity and
separation were proved from $P$ and $Z$; Section~\ref{sec:reflection}
supplies every quantifier-free induction instance. All logical
connectives, quantifiers, equality, and falsity are unchanged.
This validates every axiom and inference rule of full first-order
$\ipra$ in the target.

\subsection{Why these are bi-interpretations}

Mutual interpretability alone would be insufficient. We prove the
composites internally.

For an ordinary primitive-recursive description, use external
induction on its finite construction. Its double translation has the
same initial values and successor equation as the original, after
using the already proved equalities for earlier descriptions.
Literal-equation induction comparing the two outputs proves their
equality at every counter, with passive parameters fixed.

For a $J$ description, the double translation has initial value one
and the same increment/load transition. Induct on the single equation
equating its output with the original $J$ output. In the step,
the induction equation and the already established translation
equalities for the body identify the two instruction values.
Split coverage for that common value. The zero branch increments
equal old registers; the successor branch loads the same number.
The output equation follows in both branches. This proof uses no
induction on a formula with a quantifier.

The domains in both interpretations are defined by $x=x$.
Interpreted equality is $x=y$. The proved operation equalities show
that the graph $x=y$ is an internal total bijection preserving every
fixed symbol in both composites. Structural induction on formulas
then gives the formula round trips, since the logical operations
were left unchanged. This proves Theorem~\ref{thm:main}.
Translations preserve $\bot$ and transform finite proofs effectively,
so the equiconsistency claim follows as well.

\section{Executable artifacts and fair complexity accounting}\label{sec:artifacts}

The archive distinguishes three kinds of code.
\begin{enumerate}
\item \file{code/checker.py} checks proofs in the new
counter-aware $\TJ$ theory, including logical rules, scope conditions,
substitution, arithmetic schemas, equation induction, and the public
JSON interface. Its tests include valid proofs and invalid scope and
induction attempts.
\item The semantic evaluators and finite-description compiler execute
the displayed constructions on standard natural numbers. They are
useful implementation checks; they are not proofs of the universal
metatheorems or proof checkers.
\item \file{ancillary/checker_288.py} is a smaller AST profile of the
\emph{old} native scalar calculus. It does not validate $J$ proofs.
Its parent source and test certificates are frozen in the archive.
\end{enumerate}

All source measurements count the entire standalone module after
CPython 3.12.14 \file{ast.unparse} normalization plus one final
newline. Tokens exclude encoding, end-marker, layout, and comment
tokens; AST nodes are counted by \file{ast.walk}. Gzip uses level nine
and a zero timestamp, and is informational rather than an optimization
target. The exact machine-readable measurements are supplied in the
archive; the following table reports the relevant complete checkers.

% METRICS_TABLE_BEGIN
\begin{center}\small
\begin{tabular}{@{}lrrrr@{}}
\toprule
Complete standalone program & Bytes & Tokens & AST & Gzip-9\\
\midrule
Old byte leader &529&266&310&333\\
Old AST leader &613&263&289&349\\
Direct parent of ancillary profile &576&261&292&345\\
New ancillary old-calculus profile &604&261&288&355\\
New counter-aware $J$ reference checker &13633&2814&3060&3564\\
\bottomrule
\end{tabular}
\end{center}
% METRICS_TABLE_END

The ancillary profile replaces a recursive continuation by a loop and
uses immutable tuple contexts. It preserves the parent's prevalidated
elimination protocol and global fresh-name order. It reduces the
inspected AST minimum by one node and ties the inspected token minimum;
it is larger in bytes and gzip size. This is an implementation result
with inherited mathematics, not evidence that the new register theory
has a smaller checker.

The generic iterable opener in that ancillary profile inherits a
malformed-input qualification: in an ideal unlimited interpreter,
certain nonempty strings can cause unbounded traversal. The soundness
claim concerns finite successful executions, and completeness uses
finite canonical array/integer certificates. The new readable $J$
checker validates its concrete syntax explicitly. Neither program
assumes that a fixed physical computer checks arbitrarily large proofs.

The written metatheorems have been independently reviewed in this
research run but have not been formalized in Lean or Rocq. Finite tests
are not substitutes for those metatheorems. The delivered comparison
does not establish the requested absolute minimum, nor a new complete
checker below the old 529-byte benchmark. The positive result is the
smaller arithmetic grammar and the direct, short compiler, with the
remaining implementation cost made explicit.

\section{Further research questions}

\begin{enumerate}
\item Can a native $Q/D$ checker share its already verified equality
major premise with increment/load computation so that the new grammar
also reduces the complete trusted source? All conditional-premise,
scope, and syntax checks must be charged.
\item Is there a substantially shorter numerical selector for the
register basis? The current $C$ has a short proof but expands to a
large finite term. Source brevity, term size, and certificate size
are different objectives.
\item Does the counter-aware presentation, or the one-variable scalar
variant in Appendix~\ref{sec:scalar}, yield the smaller complete checker
after comparable implementations? The present results do not decide
this tradeoff.
\item Can the existing Lean development be reused by proving the
recursion compiler and both round trips as translations of finite
descriptions? A formalization should preserve the distinction between
fixed syntactic bodies and arbitrary functions in the metatheory.
\item Which parts of the graph specification can be weakened while
retaining actual equality and whole-domain formula round trips?
One-sided operational membership and standard-model agreement alone
are insufficient.
\end{enumerate}

\appendix
\section{A single-body-variable variant}\label{sec:scalar}

If preserving the old binding arity is preferred, use
\[
 t::=x\mid0\mid J_s(n;r.h).
\]
The register still starts at one, with the same increment/load
instructions. Remove the old counter argument from the instruction
in STEP. Define $S(a)=J_s(a;r.0)$. The $f,Z,W,D,P,G,C$ bootstrap
and intrinsic-coverage argument are unchanged because all their
bodies are constant.

For an arbitrary initialized scalar iteration with body $h(r)$, put
\begin{align*}
 H(v)&=S^3\bigl(C(Pv,a,h(P^2v))\bigr),\\
 K(n)&=J_s(Sn;v.H(v)),\\
 I(n,a;r.h)&=P^2K(n).
\end{align*}
One load computation gives, at every counter,
\[
 J_s(Sn;H)=S^2\bigl(C(P(J_s(n;H)),a,h(P^2J_s(n;H)))\bigr).
\]
Every post-initial state is therefore visibly a double successor.
At the first step, the previous register is one and $P1=0$, so
$K0=S^2a$. At the next step $PK(n)$ is a successor, and the selector
chooses the update branch. Thus
\[
 K(Sn)=S^2h(P^2K(n)),
\]
and applying $P$ twice gives the two scalar-iterator equations.
This compiler needs no induction or coverage: the displayed shape
of the state is a direct computation consequence, not an inferred
invariant.

Conversely, using the old scalar iterator, define
\[
 J_s(n;h)=I(n,1;v.C(h(v),Sv,P(h(v)))).
\]
Coverage validates its graph. Equation induction proves that its
successor abbreviation agrees with the old successor. Recurrence
uniqueness proves both composites internally, as in the main text.
Hence this scalar alternative is bi-interpretable with $\TI$ and
therefore with $\ipra$, using the repository's ordinary scalar bridge
\cite{bridge}. Its predecessor, selector, and initialization proofs
are explicitly supplied here; the existing scalar-to-general-PR
compiler is needed only for this optional variant, not for the main
counter-aware theorem.

\section{Compatibility with the inherited native logical core}

The new arithmetic can also be combined with the repository's native
choice-$D$ rules rather than ordinary logical notation. Conditional
computation needs actual proof premises. Even though J0 no longer
returns an arbitrary initializer, numerical reflexivity can still be
derived without adding a primitive reflexivity rule:
\begin{enumerate}
\item J0 for the zero body is literally $1=1$.
\item A load step with constant instruction $1=S0$ yields
$J(1;1)=0$. Transporting that equation along itself yields $0=0$.
\item Increment with the zero body and any counter $t$ yields the
literal equation $S(St)=S(St)$, since $J(t;0)=St$ by definition.
\item A load of the constant instruction $S(St)$ yields an equation
with right side $St$; self-transport yields $St=St$.
\item A load of the constant instruction $St$ yields an equation
with right side $t$; self-transport yields $t=t$.
\end{enumerate}
The derived selector then supports the inherited logical macros.
This shows that the new basis is compatible with that equality
bootstrap. It does not claim an unimplemented native $J$ checker
of the old checker's size. The executable $J$ checker delivered here
uses the explicit ordinary logical rules of Section~\ref{sec:logic}.

\begin{thebibliography}{9}
\bibitem{repo}
Vladimir Reshetnikov, \emph{ProveIt: Concise presentations of
intuitionistic primitive recursive arithmetic}, README, pinned commit
\texttt{8662bc0a00f4e77e482f7d141d7a4f919f70f9aa}.
\href{https://github.com/VladimirReshetnikov/ProveIt/blob/8662bc0a00f4e77e482f7d141d7a4f919f70f9aa/Logic/PrimitiveRecursiveArithmetic/ConciseIPRA/README.md}{Pinned README}.
\bibitem{basis}
ProveIt, \emph{Complementary introductions with one local assumption},
\file{Research/certified_zero/reference/inherited_choice_basis.tex},
at the same pinned commit.
\href{https://github.com/VladimirReshetnikov/ProveIt/blob/8662bc0a00f4e77e482f7d141d7a4f919f70f9aa/Logic/PrimitiveRecursiveArithmetic/ConciseIPRA/Research/certified_zero/reference/inherited_choice_basis.tex}{Pinned source}.
\bibitem{reflection}
ProveIt, \emph{Equality reflection and quantifier-free induction},
\file{Research/certified_zero/reference/equality_reflection.tex},
at the same pinned commit. The residual argument in
Section~\ref{sec:reflection} is an attributed adaptation.
\href{https://github.com/VladimirReshetnikov/ProveIt/blob/8662bc0a00f4e77e482f7d141d7a4f919f70f9aa/Logic/PrimitiveRecursiveArithmetic/ConciseIPRA/Research/certified_zero/reference/equality_reflection.tex}{Pinned source}.
\bibitem{bridge}
ProveIt, \emph{Arithmetic strength and both ordinary interpretation
composites; pairing and primitive recursion},
\file{Research/certified_zero/reference/inherited_ordinary_bridge.tex}
and \file{primitive_recursion.tex}, at the same pinned commit.
\href{https://github.com/VladimirReshetnikov/ProveIt/tree/8662bc0a00f4e77e482f7d141d7a4f919f70f9aa/Logic/PrimitiveRecursiveArithmetic/ConciseIPRA/Research/certified_zero/reference}{Pinned source directory}.
\bibitem{barrier}
ProveIt, \emph{Local scopes and relational obstructions},
\file{Research/local_scopes/theory/relational_obstruction.tex},
at the same pinned commit.
\href{https://github.com/VladimirReshetnikov/ProveIt/blob/8662bc0a00f4e77e482f7d141d7a4f919f70f9aa/Logic/PrimitiveRecursiveArithmetic/ConciseIPRA/Research/local_scopes/theory/relational_obstruction.tex}{Pinned source}.
\bibitem{kernels}
ProveIt, the complete checkers in \file{Research/neutral_scope},
\file{Research/coverage_and_kernel}, and
\file{Research/prevalidated_elimination}, at the same pinned commit.
Exact source paths, hashes, and reproductions are supplied with the archive.
\end{thebibliography}
\end{document}
